% Generated draft: requires mathematical and source review.
\documentclass[11pt]{article}
\usepackage[margin=25mm]{geometry}
\usepackage{fontspec}
\setmainfont{Latin Modern Roman}
\usepackage{amsmath,amssymb,mathtools,unicode-math}
\setmathfont{Latin Modern Math}
\usepackage{enumitem,needspace,xurl,hyperref}
\setlength{\parindent}{0pt}
\setlength{\parskip}{5pt}
\setlength{\emergencystretch}{3em}
\newcommand{\N}{\mathbb{N}}
\newcommand{\Z}{\mathbb{Z}}
\newcommand{\Q}{\mathbb{Q}}
\newcommand{\R}{\mathbb{R}}
\newcommand{\C}{\mathbb{C}}
\newcommand{\F}{\mathbb{F}}
\newcommand{\A}{\mathbb{A}}
\newcommand{\PP}{\mathbb{P}}
\newcommand{\E}{\mathbb{E}}
\newcommand{\eps}{\varepsilon}
\newcommand{\dd}{\,\mathrm{d}}
\begin{document}
\section*{UnsolvedMath ID: 30000014}
\textbf{Source problem code:} catalogue entry
\subsection{Definitions and mathematical statement}
\textbf{Question.} To what extent does the Floquet exponent of a Schr\"odinger operator characterize its invariant probability measure? The record supplies neither an operator model nor a definition of its Floquet exponent or invariant measure. In particular, it does not say whether the exponent is known at one spectral parameter or as an entire spectral-parameter function, or whether it means a real growth exponent or a complex Floquet quantity. These distinctions are essential to any injectivity statement.\par A conditional way to organize the question is to fix a parameter space \(\Omega\), its dynamics \(\tau\) (a transformation or flow), a specified associated family of Schr\"odinger operators \(H_\omega\), and a class \(\mathcal M_\tau(\Omega)\) of invariant probability measures. Invariance means \(\mu(\tau^{-1}B)=\mu(B)\) for a transformation, or the corresponding equality at every flow time. Once the source's exponent construction and spectral domain \(D\) are fixed, denote the resulting data by \(\mathcal F_\mu:D\to\C\) (or its specified real-valued version). Characterizing the measure from these data asks for a description of the fibers\[\{\nu\in\mathcal M_\tau(\Omega):\mathcal F_\nu(z)=\mathcal F_\mu(z)\text{ for every }z\in D\},\] and, in particular, hypotheses under which such a fiber consists of \(\mu\) alone or of explicitly identified equivalent measures. This is a framework for the missing choices, not an assertion that the excerpt has fixed \(\Omega,\tau,D\), ergodicity, or an exponent construction. No periodic, continuum, discrete, or random operator model is selected silently.\par The supplied literature triage dated 2026-08-21 labels the question open. It says that a general injectivity theorem for the Floquet exponent was not located and that later work treats special invariant measures without resolving the stated general characterization. This is a dated report of a search, not a verified nonexistence theorem or an independently established current status.
\subsection{Short English statement}
Determine what information a Schr\"odinger operator's Floquet exponent retains about the underlying invariant probability measure, and when it determines that measure uniquely. The operator family, dynamics, definition of exponent, and range of spectral data are missing from the excerpt.
\subsection{Sources}
\begin{itemize}\item Local UnsolvedMath registry, permanent ID 30000014, code catalogue entry; \emph{Floquet Exponents and Invariant Measures of Schr\"odinger Operators}, Oberwolfach Reports Open Problems.\item \emph{Mathematics and Physics of Disordered Systems (2005)}, Oberwolfach Reports, DOI 10.4171/owr/2004/22, pp. 1167--1232. \url{https://doi.org/10.4171/owr/2004/22}. Proposed year and year tag: 2004; no proposer named.\item Supplied literature source labeled \emph{official report}: \url{https://ems.press/content/serial-article-files/45912?nt=1}.\item Literature triage checked 2026-08-21, status \texttt{open}; record updated 2026-09-16. The \texttt{corrected\_verified} statement label concerns spelling of Schr\"odinger, not recovery of missing operator hypotheses. Verify the dated status and claims independently.\end{itemize}
\end{document}
